\clearpage

\section{ALCANCE}

    El presente proyecto abarca el análisis, diseño, desarrollo e implementación de
    una plataforma web funcional orientada a gestionar el ciclo operativo del Desafío
    Bebras en Bolivia, con foco en estudiantes de nivel primario y secundario.
    El alcance del proyecto se estructura en los siguientes componentes:

    \begin{itemize}
        \item Definición de políticas y reglas de negocio que regulen el funcionamiento del concurso en el contexto boliviano, en correspondencia con los lineamientos internacionales de Bebras, constituyendo la base normativa sobre la cual operará el sistema.

        \item Desarrollo de una plataforma web funcional basada en una arquitectura cliente-servidor, organizada en módulos que representan las principales actividades del ciclo del concurso:
        \begin{itemize}
            \item Gestión de competencias y tareas de pensamiento computacional.
            \item Gestión de contenido institucional y sitio web informativo.
            \item Registro y administración de unidades educativas escolares, docentes y estudiantes de nivel primario y secundario.
            \item Mecanismos de evaluación digital e imprimible para entornos escolares.
        \end{itemize}

        \item Validación con usuarios mediante pruebas funcionales que permitan evaluar la contribución de la plataforma a la organización y ejecución del Desafío Bebras Bolivia en el ámbito escolar.
    \end{itemize}

\section{LÍMITES}

    El proyecto se delimita bajo las siguientes restricciones:

    \begin{itemize}
        \item La plataforma se orienta exclusivamente a la gestión del Desafío Bebras en Bolivia para estudiantes de nivel primario y secundario, sin extenderse a otros niveles del sistema educativo como la educación superior o la formación técnica.

        \item El sistema desarrollado corresponde a una solución funcional en entorno controlado, por lo que no contempla características propias de plataformas a escala institucional amplia, como alta concurrencia masiva, seguridad avanzada de nivel corporativo o integración con múltiples servicios externos.

        \item No se incluyen módulos relacionados con autenticación institucional gubernamental, gestión curricular escolar, comunicación directa con sistemas del Ministerio de Educación ni integración técnica con plataformas educativas nacionales existentes.

        \item El despliegue del sistema se limita a un entorno controlado de desarrollo y pruebas, sin contemplar su puesta en producción a escala nacional como parte del presente proyecto.

        \item La validación del sistema se realizará con el equipo organizador y una muestra representativa de estudiantes participantes, sin contemplar la ejecución completa de una edición nacional del concurso como parte del proyecto.
    \end{itemize}
